\capitulo{3}{Conceptos teóricos}

En aquellos proyectos que necesiten para su comprensión y desarrollo de unos conceptos teóricos de una determinada materia o de un determinado dominio de conocimiento, debe existir un apartado que sintetice dichos conceptos.

Algunos conceptos teóricos de \LaTeX{} \footnote{Créditos a los proyectos de Álvaro López Cantero: Configurador de Presupuestos y Roberto Izquierdo Amo: PLQuiz}.

\section{Sistemas SIEM}

Un sistema SIEM (Security Information and Event Management) es una plataforma diseñada para recopilar, almacenar, analizar y correlacionar eventos de seguridad procedentes de diferentes fuentes dentro de una organización.

Su principal objetivo es proporcionar una visión centralizada de la actividad del sistema, facilitando la detección de incidentes de seguridad, comportamientos anómalos y posibles amenazas. Para ello, los sistemas SIEM reciben información procedente de dispositivos, aplicaciones, sistemas operativos, sensores y otros elementos de la infraestructura monitorizada.

Entre las funcionalidades más habituales de los sistemas SIEM se encuentran la recopilación de logs, la correlación de eventos, la generación de alertas y la monitorización en tiempo real de los distintos activos supervisados.

TODO
Añadir esquema tipo así Dispositivos logs -siem- correlación-alertas-dashboard


\section{Gestión de eventos y logs}

Los logs constituyen la principal fuente de información utilizada por un sistema SIEM. Un log es un registro generado por un dispositivo o aplicación que contiene información sobre una acción, evento o incidencia ocurrida en un momento determinado.

La correcta estructuración de los eventos resulta fundamental para facilitar su almacenamiento, consulta y posterior análisis. En este proyecto se utilizan eventos estructurados que contienen información relacionada con el origen del evento, su severidad, el dispositivo que lo genera y una descripción detallada de la situación detectada.

\section{Correlación de eventos}

La correlación de eventos consiste en analizar conjuntamente varios registros con el objetivo de detectar patrones o situaciones que podrían pasar desapercibidas si se analizaran de forma individual.

Mediante la aplicación de reglas de correlación es posible identificar comportamientos sospechosos, accesos no autorizados o combinaciones de eventos que indiquen la existencia de un posible incidente de seguridad.

\section{Detección de anomalías}

La detección de anomalías permite identificar comportamientos que se desvían del funcionamiento esperado del sistema. Este tipo de análisis resulta especialmente útil para detectar incidencias operativas o situaciones potencialmente peligrosas.

En el proyecto desarrollado se han implementado mecanismos básicos de detección de anomalías relacionados con accesos fuera de horario, accesos denegados repetidos, condiciones ambientales fuera de rango y actividad anómala en determinados depósitos.

\section{Infraestructuras críticas}

Las infraestructuras críticas son instalaciones cuya interrupción o funcionamiento incorrecto puede generar consecuencias importantes sobre la seguridad, la economía o el funcionamiento de servicios esenciales.

Debido a la importancia de estas instalaciones, resulta necesario disponer de mecanismos de monitorización capaces de supervisar continuamente su estado y generar alertas ante cualquier situación que pueda comprometer su seguridad o disponibilidad.

\subsection{Subsecciones}

Además de secciones tenemos subsecciones.

\subsubsection{Subsubsecciones}

Y subsecciones. 


\section{Referencias}

Las referencias se incluyen en el texto usando cite~\cite{wiki:latex}. Para citar webs, artículos o libros~\cite{koza92}, si se desean citar más de uno en el mismo lugar~\cite{bortolot2005, koza92}.


\section{Imágenes}

Se pueden incluir imágenes con los comandos standard de \LaTeX, pero esta plantilla dispone de comandos propios como por ejemplo el siguiente:

\imagen{escudoInfor}{Autómata para una expresión vacía}{.5}



\section{Listas de items}

Existen tres posibilidades:

\begin{itemize}
	\item primer item.
	\item segundo item.
\end{itemize}

\begin{enumerate}
	\item primer item.
	\item segundo item.
\end{enumerate}

\begin{description}
	\item[Primer item] más información sobre el primer item.
	\item[Segundo item] más información sobre el segundo item.
\end{description}
	
\begin{itemize}
\item 
\end{itemize}

\section{Tablas}

Igualmente se pueden usar los comandos específicos de \LaTeX o bien usar alguno de los comandos de la plantilla.

\tablaSmall{Herramientas y tecnologías utilizadas en cada parte del proyecto}{l c c c c}{herramientasportipodeuso}
{ \multicolumn{1}{l}{Herramientas} & App AngularJS & API REST & BD & Memoria \\}{ 
HTML5 & X & & &\\
CSS3 & X & & &\\
BOOTSTRAP & X & & &\\
JavaScript & X & & &\\
AngularJS & X & & &\\
Bower & X & & &\\
PHP & & X & &\\
Karma + Jasmine & X & & &\\
Slim framework & & X & &\\
Idiorm & & X & &\\
Composer & & X & &\\
JSON & X & X & &\\
PhpStorm & X & X & &\\
MySQL & & & X &\\
PhpMyAdmin & & & X &\\
Git + BitBucket & X & X & X & X\\
Mik\TeX{} & & & & X\\
\TeX{}Maker & & & & X\\
Astah & & & & X\\
Balsamiq Mockups & X & & &\\
VersionOne & X & X & X & X\\
} 
